\section{The integral local factors}\label{sec:intro}

The universal count of effective curves on a compact threefold records
the entire space of cycles in each integral homology class. Even when
these spaces carry integral classes, the factors used to express their
generating series can have fractional exponents. On a smooth projective
Calabi--Yau threefold, for example, a genus-ten factor occurs with
exponent in $\mathbb Z+2/3$ in three times a primitive class
\cite[Theorem~1.1]{WueCounter}, \cite[Theorem~1.2]{WueTriples}. This factor comes from scaling the two
normal directions to a curve by the same weight. Which local factors
recover the integral structure?

Pardon constructed an algebra of effective one-cycles in families of
complex threefolds. In its semi-Fano part every irreducible component
has nonnegative first Chern number. Write $H_{\ZZ}$ for the summand
in first Chern number and virtual dimension zero. Multiplication is
external product on disjoint targets. Pardon proved that integral
geometric elements obtained by isolating multiples of embedded curves,
one for each genus and multiplicity, freely generate this algebra.
The normal neighborhood used to isolate a multiple may depend on its
multiplicity. A uniform construction by scalar dilation gives free
generators over $\QQ$ \cite[Theorem~1.2]{Pardon}. We identify a uniform
equivariant construction which gives free generators over $\ZZ$.

Let $C$ be a smooth compact curve of genus $g$. Choose a line bundle
$L$ of degree $g-1$ and give the two summands of
$L\oplus(K_C\otimes L^{-1})$ weights $+1$ and $-1$. We call this
the balanced antidiagonal model. Write $a_{g,m}$ for the normalized
equivariant class supported on the $m$-fold zero section and
$A_g(q)=1+\sum_{m\ge1}a_{g,m}q^m$. This class is independent of the
choices of $C$ and $L$. The normalization divides a supported Thom
class by a product of weights, so its construction gives an element of
$H=H_{\ZZ}\otimes\QQ$. For equal weights write $x_{g,m}$ for the scalar
class. The two classes agree in multiplicity one.

\begin{theorem}\label{thm:basis}
The map to the rationalization embeds $H_{\ZZ}$ as the polynomial ring
$\ZZ[a_{g,m}:g\ge0,\ m\ge1]$. For each $(g,m)$ one can choose an
integral geometric local curve element $y_{g,m}$ of this type so that
\[
 y_{g,m}-a_{g,m}\ \in\ \ZZ[a_{h,j}:h\ge0,\ 1\le j<m].
\]
\end{theorem}

There is a geometric obstruction to comparing the two sorts of
generators directly in the normal bundle. For $g\ge2$ choose $L$ outside
the theta divisor. Both normal summands are acyclic, so the embedded
zero section is regular and isolated. Yet $h^0(C,L^2)\ge g-1$.
A nonzero section $s$ cuts out the compact double cycles $w^2=ts$ in
the first summand. They converge to the doubled zero section as $t$
tends to zero. Isolation of the embedded curve says little about
isolation of its doubles.

For a prescribed multiplicity $m$ we start in a holomorphic
neighborhood in which all multiples through $m$ are isolated. We place a tube of this neighborhood in a closed symplectic six-torus.
Bounded-area compactness keeps the cycles in the tube away from its
boundary during a small deformation. Excision therefore lets us apply
the compact cluster comparison to its insertion, giving a product of contributions of elementary curves, the regular local
models of Ionel and Parker \cite{IP}. A genus-$h$ elementary curve
contributes the balanced series $A_h$ by the geometric comparison of
\cite[Proposition~3.1]{WueProduct}. In degree one, testing by genus
separates the multiplicity-one classes: their fixed-genus series have
distinct leading powers. The embedded curve has its complex orientation,
so its genus occurs with coefficient $+1$ and every other genus with
coefficient zero. The remaining factors have core
class a proper multiple of the original class. In the coefficient of
$q^m$ these corrections therefore involve only generators of
multiplicity less than $m$. This gives the integral triangular
comparison in Theorem~\ref{thm:basis}.

For a smooth compact Kähler threefold $X$, write $\Gamma_0(X)$ for zero
and the effective integral classes of first Chern number zero. Let
$U_{X,\beta}$ be the characteristic function of the entire space of
effective cycles of class $\beta$, and take $U_{X,0}=1$.
\begin{theorem}\label{thm:compact}
If $c_1(TX)$ is nonnegative on every compact irreducible curve, then
\begin{equation}\label{eq:compact}
 \sum_{\beta\in\Gamma_0(X)}U_{X,\beta}Q^\beta
 =\prod_{\substack{\beta\in\Gamma_0(X)\\ \beta\ne0}}\prod_{h\ge0}
           A_h(Q^\beta)^{n_{\beta,h}}
 \quad\text{in }H_{\ZZ}[[\Gamma_0(X)]] .
\end{equation}
The exponents are unique integers and only finitely many occur for each
$\beta$. The completion is by Kähler area, and the integers are independent
of the Kähler form. When $X$ is smooth projective with trivial canonical
bundle and $H^1(X,\cO_X)=H^2(X,\cO_X)=0$, the exponent $n_{\beta,h}$
is the signed multiplicity of elementary clusters of class $\beta$ and
genus $h$ in the bounded-energy cluster decomposition.
\end{theorem}

Two identities for cycles explain why the basis theorem yields a
product. Addition records every decomposition of a cycle, and division
by an integer records every integral root of its class. After the
lower-area factors have been removed, the first remaining coefficient
is primitive for addition and vanishes under every proper division.
It is therefore an integral combination of multiplicity-one classes.
Their coefficients give the exponents at that homology label.
The product retains the full integral homology label, including
torsion. Its local factors also have a numerical interpretation.
Bryan and Pandharipande computed the local Gromov--Witten residues
\cite{BP}, and Ionel and Parker identified the antidiagonal residue of an
elementary curve \cite[Proposition~3.3]{IP}. Doan, Ionel and Walpuski
proved the requisite finiteness for the symplectic cluster expansion
\cite{DIW}. The balanced universal classes and a conditional numerical
factorization were constructed in \cite{WueBalanced}. The geometric
comparison of \cite{WueProduct} established the projective cluster
formula. Here its use on a tube gives the integral presentation and
the compact Kähler factorization. In the compact example above, the
primitive genus-three count is $-14$. Changing back to scalar factors
subtracts $(4720/3)(-14)$ from the genus-ten exponent, giving the stated
fractional part \cite{WueCounter,WueTriples}.

\subsection*{Acknowledgements}
I thank John Pardon for pointing out the difficulty of comparing his
scalar-equivariant local curve classes with the elementary curves of
Ionel and Parker.
